\documentclass[10pt,a4paper]{amsart}
\usepackage{amsmath,amssymb,amsthm,graphicx}
% Retained preamble for the Stacks Project chapter snapshot in this directory.
% The chapter's own preamble.tex is a copy of the shared Stacks Project preamble; the block
% below is sliced MECHANICALLY, line for line, from it. Nothing was added, removed or
% redefined; the wrapper class, the xr-hyper/hyperref wiring and the externaldocument list
% are omitted because the assembler supplies the document class.
%
% lines taken from preamble.tex: 3, 5, 9, 11, 12, 156, 157, 158, ... (51 lines total)
%
% package coverage restored from the chapter's own preamble (2026-10-13): xy, xy-2cell

\usepackage[all]{xy}
\xyoption{2cell}
\UseAllTwocells

\newenvironment{reference}{\comment}{\endcomment}
\newenvironment{slogan}{\comment}{\endcomment}
\newenvironment{history}{\comment}{\endcomment}
\theoremstyle{plain}
\newtheorem{theorem}[subsection]{Theorem}
\newtheorem{proposition}[subsection]{Proposition}
\newtheorem{lemma}[subsection]{Lemma}
\theoremstyle{definition}
\newtheorem{definition}[subsection]{Definition}
\newtheorem{example}[subsection]{Example}
\newtheorem{exercise}[subsection]{Exercise}
\newtheorem{situation}[subsection]{Situation}
\theoremstyle{remark}
\newtheorem{remark}[subsection]{Remark}
\newtheorem{remarks}[subsection]{Remarks}
\numberwithin{equation}{subsection}
\def\lim{\mathop{\mathrm{lim}}\nolimits}
\def\colim{\mathop{\mathrm{colim}}\nolimits}
\def\Spec{\mathop{\mathrm{Spec}}}
\def\Hom{\mathop{\mathrm{Hom}}\nolimits}
\def\Ext{\mathop{\mathrm{Ext}}\nolimits}
\def\SheafHom{\mathop{\mathcal{H}\!\mathit{om}}\nolimits}
\def\SheafExt{\mathop{\mathcal{E}\!\mathit{xt}}\nolimits}
\def\Sch{\mathit{Sch}}
\def\Mor{\mathop{\mathrm{Mor}}\nolimits}
\def\Ob{\mathop{\mathrm{Ob}}\nolimits}
\def\Sh{\mathop{\mathit{Sh}}\nolimits}
\def\NL{\mathop{N\!L}\nolimits}
\def\CH{\mathop{\mathrm{CH}}\nolimits}
\def\proetale{{pro\text{-}\acute{e}tale}}
\def\etale{{\acute{e}tale}}
\def\QCoh{\mathit{QCoh}}
\def\Ker{\mathop{\mathrm{Ker}}}
\def\Im{\mathop{\mathrm{Im}}}
\def\Coker{\mathop{\mathrm{Coker}}}
\def\Coim{\mathop{\mathrm{Coim}}}
\DeclareMathSymbol{\boxtimes}{\mathbin}{AMSa}{"02}
\def\QCohstack{\mathcal{QC}\!\mathit{oh}}
\def\Cohstack{\mathcal{C}\!\mathit{oh}}
\def\Spacesstack{\mathcal{S}\!\mathit{paces}}
\def\Quotfunctor{\mathrm{Quot}}
\def\Hilbfunctor{\mathrm{Hilb}}
\def\Curvesstack{\mathcal{C}\!\mathit{urves}}
\def\Polarizedstack{\mathcal{P}\!\mathit{olarized}}
\def\Complexesstack{\mathcal{C}\!\mathit{omplexes}}
\def\Pic{\mathop{\mathrm{Pic}}\nolimits}
\def\Picardstack{\mathcal{P}\!\mathit{ic}}
\def\Picardfunctor{\mathrm{Pic}}
\def\Deformationcategory{\mathcal{D}\!\mathit{ef}}
\usepackage{cleveref}
\crefname{originalequation}{original equation}{original equations}
\Crefname{originalequation}{Original equation}{Original equations}
\makeatletter
\newlabel{equation-fibre-product}{{09DK}{source}{}{sourcepointer-1}{}}
\newlabel{equation-fibre-product@cref}{{[originalequation][0][]09DK}{[1][0][]source}{}{}{}}
\makeatother
\crefname{originallemma}{original lemma}{original lemmas}
\Crefname{originallemma}{Original lemma}{Original lemmas}
\makeatletter
\newlabel{lemma-pullback-cotangent-morphism-topoi}{{08SV}{source}{}{sourcepointer-2}{}}
\newlabel{lemma-pullback-cotangent-morphism-topoi@cref}{{[originallemma][0][]08SV}{[1][0][]source}{}{}{}}
\makeatother
\crefname{originallemma}{original lemma}{original lemmas}
\Crefname{originallemma}{Original lemma}{Original lemmas}
\makeatletter
\newlabel{lemma-stalk-cotangent-complex}{{08T0}{source}{}{sourcepointer-3}{}}
\newlabel{lemma-stalk-cotangent-complex@cref}{{[originallemma][0][]08T0}{[1][0][]source}{}{}{}}
\makeatother
\crefname{originallemma}{original lemma}{original lemmas}
\Crefname{originallemma}{Original lemma}{Original lemmas}
\makeatletter
\newlabel{lemma-tensor-product}{{09DB}{source}{}{sourcepointer-4}{}}
\newlabel{lemma-tensor-product@cref}{{[originallemma][0][]09DB}{[1][0][]source}{}{}{}}
\makeatother
\crefname{originalsection}{original section}{original sections}
\Crefname{originalsection}{Original section}{Original sections}
\makeatletter
\newlabel{section-localization}{{08QY}{source}{}{sourcepointer-5}{}}
\newlabel{section-localization@cref}{{[originalsection][0][]08QY}{[1][0][]source}{}{}{}}
\makeatother
\title{The cohomology of the cone attached to a fibre product of algebraic spaces}
\author{The Stacks Project Authors}
\date{stacks-project:cotangent.tex:local-snapshot-20261012; GNU FDL 1.2 or later}
\begin{document}
\maketitle
\setcounter{subsection}{1}
\noindent\textit{Editorial note.} All delivered mathematical text is the author's own, copied from the pinned Stacks Project chapter file cotangent.tex; no statement or proof step was rewritten, reordered or removed.\par\medskip
\noindent\textit{Editorial note.} The statement displays the degree -2 identification and fixes the four local rings in the same display; the proof is a stalk reduction followed by the ring case.\par\medskip
\noindent\textit{Editorial note.} The five references of the counted spans are recorded as non-load-bearing navigation pointers; the argument that consumes them stays inside the counted proof span.\par\medskip
\noindent\textit{Editorial note.} Preamble: the retained preamble of this chapter (sources/stacks-cotangent/retained-preamble.tex, lines 1-63), containing the Xy-pic setup, the comment environments, the theorem declarations and the macros including \textbackslash{}NL.\par\medskip
\noindent\textit{Allowed standard prerequisites (editorial).}\begin{itemize}
\item Commutation of the cotangent complex with stalks and with pullback. Lemmas lemma-stalk-cotangent-complex and lemma-pullback-cotangent-morphism-topoi, used at proof lines 4841-4844. Reduces the statement to a computation at one geometric point.
\item The localization results for the cotangent complex. Section section-localization, used at proof lines 4845-4847. Give the identification L\_\{C/R\} = L\_\{A tensor\_R B/R\} tensor C.
\item The henselization description of the etale local ring of a fibre product. used at proof lines 4845-4845. Identifies C as a henselization of A tensor\_R B.
\item The ring case of the computation. Lemma lemma-tensor-product (3063-3076, produced as L44133's blocked counterpart; here reported as the ring-level input), used at proof lines 4850-4851. Supplies the vanishing and the degree -2 identification.
\end{itemize}\medskip
\ifdefined\hypertarget\hypertarget{sourcepointer-1}{}\fi
\noindent\textit{Source reference.} Original equation 09DK: focused-seeds/AGP/sources/stacks-snapshot-20261011/tags.txt, source lines 11868--11868. Navigation pointer into the Stacks Project; the standard fact it supplies is applied inside the counted proof span and the argument resting on it stays there.\par\medskip
\ifdefined\hypertarget\hypertarget{sourcepointer-2}{}\fi
\noindent\textit{Source reference.} Original lemma 08SV: focused-seeds/AGP/sources/stacks-snapshot-20261011/tags.txt, source lines 11143--11143. Navigation pointer into the Stacks Project; the standard fact it supplies is applied inside the counted proof span and the argument resting on it stays there.\par\medskip
\ifdefined\hypertarget\hypertarget{sourcepointer-3}{}\fi
\noindent\textit{Source reference.} Original lemma 08T0: focused-seeds/AGP/sources/stacks-snapshot-20261011/tags.txt, source lines 11148--11148. Navigation pointer into the Stacks Project; the standard fact it supplies is applied inside the counted proof span and the argument resting on it stays there.\par\medskip
\ifdefined\hypertarget\hypertarget{sourcepointer-4}{}\fi
\noindent\textit{Source reference.} Original lemma 09DB: focused-seeds/AGP/sources/stacks-snapshot-20261011/tags.txt, source lines 11859--11859. Navigation pointer into the Stacks Project; the standard fact it supplies is applied inside the counted proof span and the argument resting on it stays there.\par\medskip
\ifdefined\hypertarget\hypertarget{sourcepointer-5}{}\fi
\noindent\textit{Source reference.} Original section 08QY: focused-seeds/AGP/sources/stacks-snapshot-20261011/tags.txt, source lines 11076--11076. Navigation pointer into the Stacks Project; the standard fact it supplies is applied inside the counted proof span and the argument resting on it stays there.\par\medskip
\begin{lemma}
\label{lemma-fibre-product}
Let $S$ be a scheme. Let $X \to B$ and $Y \to B$ be morphisms of algebraic
spaces over $S$. The object $E$ in (\ref{equation-fibre-product}) satisfies
$H^i(E) = 0$ for $i = 0, -1$ and for a geometric point
$(\overline{x}, \overline{y}) : \Spec(k) \to X \times_B Y$ we have
$$
H^{-2}(E)_{(\overline{x}, \overline{y})} =
\text{Tor}_1^R(A, B) \otimes_{A \otimes_R B} C
$$
where $R = \mathcal{O}_{B, \overline{b}}$, $A = \mathcal{O}_{X, \overline{x}}$,
$B = \mathcal{O}_{Y, \overline{y}}$, and
$C = \mathcal{O}_{X \times_B Y, (\overline{x}, \overline{y})}$.
\end{lemma}

\begin{proof}
The formation of the cotangent complex commutes with taking stalks
and pullbacks, see
Lemmas \ref{lemma-stalk-cotangent-complex} and
\ref{lemma-pullback-cotangent-morphism-topoi}.
Note that $C$ is a henselization of $A \otimes_R B$.
$L_{C/R} = L_{A \otimes_R B/R} \otimes_{A \otimes_R B} C$
by the results of Section \ref{section-localization}.
Thus the stalk of $E$ at our geometric point is the cone of the
map $L_{A/R} \otimes C \to L_{A \otimes_R B/R} \otimes C$.
Therefore the results of the lemma follow from
the case of rings, i.e., Lemma \ref{lemma-tensor-product}.
\end{proof}

\end{document}
